% Research proposal scaffold, not a completed submission.
% Target and verification record: ../guides/SUBMISSION_PLAN.md
\documentclass[letterpaper,10pt,conference]{ieeeconf}
\overrideIEEEmargins
\usepackage{amsmath,amssymb,booktabs,graphicx,cite}
\usepackage[pdfauthor={},pdftitle={Replay-Stable Reduction of Helicopter Hover-Control Failure Scenarios}]{hyperref}
\title{\LARGE\bf Replay-Stable Reduction of Helicopter Hover-Control Failure Scenarios}
% Anonymous review: no author block, affiliation, identifying funding or repository URL.
\begin{document}
\maketitle
\thispagestyle{empty}
\pagestyle{empty}
\begin{abstract}
This proposal investigates reducing simulation scenarios that violate helicopter hover-control requirements while retaining the same violation under a declared set of replay conditions. The planned study compares event deletion and continuous-parameter shrinking with and without replay-aware acceptance under matched simulation budgets. Model validation, controller implementation and evaluation remain pending; no improvement or novelty claim is established.
\end{abstract}

\section{Introduction}
% Motivate actionable failure reproducers. State measured contributions only after evaluation.
% Planned contribution: a replay-aware reduction predicate and an empirical trade-off study.

\section{Related Work}
Time-structured falsification has been investigated previously~\cite{ernst2022}.
Prior work includes debugging a small helicopter control model~\cite{singh2020} and minimizing bug-triggering robotic-vehicle inputs~\cite{wang2026}.
Stochastic failure-preserving reduction of cyber-physical test inputs is already addressed in a 2026 preprint~\cite{valle2026}; our more specific numerical-replay hypothesis is not yet shown to be novel.
% Explain precise distinctions; a helicopter application alone does not establish novelty.

\section{Helicopter Model and Validation}
% Aircraft provenance/license, model equations/reference, controller, units/frames, trim,
% stable position hold, numerical convergence and external-validation limitations.

\section{Scenario Reduction Method}
\subsection{Admissible Disturbances and Failure Signature}
% Smooth bounded gusts, timestamped delay, fixed controller/sensor clocks and dwell semantics.
\subsection{Replay-Aware Acceptance}
% Define replay-set predicate and violation margin; no global-minimality or physical-safety claim.
\subsection{Reduction Operators and Budget}
% Event deletion then continuous shrinking; count all validation simulations.

\section{Experimental Design}
% Shared failure corpus, development/evaluation split, nominal-only and post-hoc baselines,
% matched compute, held-out replay tests, clustering unit and uncertainty intervals.

\section{Results and Discussion}
No confirmatory experimental results are available.
% Generate all tables/figures from evidence; include failed reductions and negative findings.

\section{Limitations and Threats to Validity}
The proposed study concerns one simulation testbed. Real-helicopter behaviour and generalization to other aircraft have not been established. The existing feasibility probe is not a validated fixed-position hover experiment.

\section{Conclusion}
% Write after evaluation; do not predict outcomes.

\section*{AI-Use Disclosure}
OpenAI Codex generated this proposal scaffold and assisted with research planning and preliminary code. A final disclosure must identify the affected manuscript sections and extent of use. Scientific validation and human author review remain pending.
% Review disclosure placement with conference instructions before upload; do not silently omit.

\bibliographystyle{IEEEtran}
\bibliography{references}
\end{document}
